\chapter{Conclusions and Future Work}

\section{Conclusion}
Phase 1 of this project analysed why Large Language Models struggle with long interactions: they are
stateless, their effective context is much shorter than their nominal window, and retrieval alone does
not stop hallucination. It proposed a persistent memory layer as the remedy. Phase 2 built that layer and
measured it.

The implemented system stores every conversation turn as a normalised, third-person note with an
importance score, an embedding, BM25 tokens and extracted entities. At query time it fuses dense and
lexical retrieval, expands candidates one hop through the entity graph, re-ranks them with a
cross-encoder, scores them with relevance-gated recency/importance/coherence terms, resolves conflicting
versions of a fact in favour of the newest one, injects the top five notes into a 512-token prompt and
finally verifies the generated answer against those notes.

On a \nQuestions-question long-term conversational benchmark, answered by the same 1.5B-parameter local
model under every strategy, the memory layer with its grounding gate reached \accMemGate\% accuracy with a
hallucination rate of \halMemGate\%, compared with \accFull\% and \halFull\% for sending the full
conversation history, \accRag\% and \halRag\% for the Phase-1 vector-RAG design, and \accWin\% for a
sliding window of recent turns. It did so with \tokMem{} prompt tokens per query, \tokRatioFullMem{} times
fewer than the full context, and \latRatioFullMem{} times lower latency. The stage-wise ablation showed
where the gains come from: the entity graph lifted multi-hop retrieval recall from \rFiveVecMH\% to
\rFiveMemMH\%, the conflict rule kept both versions of updated facts available, and the question penalty
removed small-talk turns that lexically resemble questions. Importance-based pruning showed that the memory
store can be halved without loss of recall, or reduced to a sixth of the turns at a cost of under four
recall points. The grounding gate, a set of
three deterministic evidence checks, detected hallucinations with precision \gatePMem{} and F1 \gateFMem,
against precision 0.25 and F1 0.39 for the detectors examined in Phase 1, at zero additional model cost.

The context-length experiment reproduced on our own hardware what Phase 1 took from the literature: the
cost of the full-context approach grows linearly with history length (\niahTokFullLast{} tokens and
\niahLatFullLast{}~s per query at \niahMaxLen{} tokens) while the memory layer stays constant
(\niahTokMemLast{} tokens, \niahLatMemLast{}~s). More importantly, the main benchmark showed that the real
difficulty of long-term memory is not length but \emph{interference} between many similar, updated and
unrelated facts, and that selecting a few verified notes removes that interference.

Taken together, the two phases support the thesis set out at the start of the project: a persistent,
selective and verified memory layer is a better foundation for long-running LLM applications than a larger
context window, and it can be built from standard, inexpensive components.

\section{Limitations}
\begin{itemize}[nosep]
  \item The benchmark is synthetic. It is controlled and reproducible, and it follows the LoCoMo and
        LongMemEval question taxonomy, but its conversations are simpler than real ones.
  \item The answering model is small (1.5B parameters). Retrieval recall is above 96\% while accuracy is
        about \accMem\%; most of the remaining errors are reading errors that a larger model would avoid.
  \item Entity extraction, importance scoring, query expansion and conflict detection are rule-based. They
        are free and transparent, but they only capture surface patterns.
  \item Hyper-parameters were tuned on a development split generated by the same generator; transfer to
        other domains would need re-tuning.
\end{itemize}

\section{Future Work}
\subsection{Write-Time Fact Extraction}
Replace the raw turn with LLM-extracted atomic facts (as in Mem0) so that a single turn containing two facts
becomes two memories, and so that the conflict rule can compare facts rather than sentences. The
extraction can be batched off-line, keeping query latency unchanged.

\subsection{Learned Components}
Fine-tune the cross-encoder on (question, memory note) pairs from the benchmark so that re-ranking is
reliable on conversational text, and train a light classifier for importance and for the
\emph{about-the-user} decision that the gate currently makes with regular expressions.

\subsection{Typed Entity Graph}
Extend the memory--entity graph with relation types (\emph{sister-of}, \emph{works-at},
\emph{lives-in}) so that multi-hop questions can be answered by typed traversal rather than one-hop
expansion, and so that an update of a relation (a move) can retire the old edge explicitly.

\subsection{Consolidation and Forgetting}
Periodically merge related notes into summaries and apply the Ebbinghaus-style decay of MemoryBank to
notes that are never accessed, keeping the store bounded for conversations that run for years.

\subsection{Evaluation on Public Benchmarks and Larger Models}
Run the layer on LoCoMo and LongMemEval with 7B--70B models to confirm that the gains transfer, and
measure the trade-off between answering-model size and memory-layer sophistication.

\subsection{Deployment}
Package the layer as a service with per-user isolation, encryption of stored notes, a right-to-forget
operation, and a web demo in which the retrieved notes and the gate decision are shown next to each
answer so that users can audit what the assistant remembers.
